% !Mode:: "TeX:UTF-8"
% !TEX root = ../../main.tex
% ============================================================
% TOPIC 24
% ============================================================
\mashgulot{24}{The Seidel method for the approximate solution of systems of linear algebraic equations}

\begin{tarif}
The \textbf{Seidel (Gauss–Seidel) method} is a modification of the simple
iteration method: when $x_i^{(k+1)}$ is computed, the values
$x_1^{(k+1)},\dots,x_{i-1}^{(k+1)}$ already found in the current iteration are
used at once:
\[\kf{x_i^{(k+1)}=\frac{1}{a_{ii}}\Big(b_i-\sum\limits_{j<i}a_{ij}x_j^{(k+1)}-\sum\limits_{j>i}a_{ij}x_j^{(k)}\Big),
\qquad i=1,\dots,n .}\]
\end{tarif}

\begin{teorema}
If $A$ is diagonally dominant, or if $A$ is symmetric and positive definite,
the Seidel method converges to the unique solution for every initial
approximation.
\end{teorema}

As a rule, the Seidel method converges faster than simple iteration, and
only one vector has to be stored during the computation.

\begin{misol}
Solve the system
$\begin{cases}10x_1+x_2+x_3=12,\\2x_1+10x_2+x_3=13,\\2x_1+2x_2+10x_3=14\end{cases}$
of the previous topic by the Seidel method.

\textbf{Solution.} We start from $x^{(0)}=(0;0;0)$. In the first iteration
each new value is substituted immediately into the next equation:
\[\kf{x_1^{(1)}=\frac{12-0-0}{10}=1.2,\quad
x_2^{(1)}=\frac{13-2\cdot1.2-0}{10}=1.06,\quad
x_3^{(1)}=\frac{14-2\cdot1.2-2\cdot1.06}{10}=0.948 .}\]
The second iteration uses $x^{(1)}$ in the same way:
\[\kf{x_1^{(2)}=\frac{12-1.06-0.948}{10}=0.9992,\quad
x_2^{(2)}=\frac{13-1.9984-0.948}{10}=1.00536,\quad
x_3^{(2)}=\frac{14-1.9984-2.01072}{10}=0.99909 .}\]
Continuing, we obtain the table below; the last column gives the error with
respect to the exact solution $x^*=(1;1;1)$.

\begin{center}
\begin{tabular}{c|ccc|c}
$k$ & $x_1^{(k)}$ & $x_2^{(k)}$ & $x_3^{(k)}$ & $\|x^{(k)}-x^*\|_\infty$\\ \hline
1 & 1.2 & 1.06 & 0.948 & 0.2\\
2 & 0.9992 & 1.00536 & 0.99909 & 0.0054\\
3 & 0.99956 & 1.00018 & 1.00005 & 0.00044\\
\end{tabular}
\end{center}
After three iterations the error of the Seidel method is about
$4.4\cdot10^{-4}$, whereas simple iteration gave $0.03$: the Seidel method
converges considerably faster here.
\end{misol}

\begin{eslatma}
The iteration is usually stopped as soon as
$\|x^{(k+1)}-x^{(k)}\|<\varepsilon$, where $\varepsilon$ is the prescribed
accuracy.
\end{eslatma}

\subsection{The Seidel method in information security}

\textbf{Devices with limited resources.} IoT devices and smart cards have
very little memory. The Seidel method overwrites the old values with the new
ones as it goes (“in-place” computation), which makes it convenient for
approximate computations in embedded cryptographic modules.

\textbf{Network models and traffic distribution.} The flow balance equations
between the nodes of a network are diagonally dominant: the load of each node
exceeds the influence of its neighbours. Redistributing the load during a DDoS
attack requires such systems to be solved again quickly, and the Seidel method
suits this well, since the previous solution is a good initial approximation.

\textbf{Iterative methods and side-channel attacks.} If the number of
iterations depends on the data, secret information may “leak” through the
running time. For this reason cryptographic software should use a number of
iterations that is fixed in advance.

\subsection{Computing with Python, Maple and Mathcad}

The Seidel iteration is easy to program, because the vector \texttt{x} is
updated in place. The following code repeats the example above and prints
the error of each approximation:
\begin{Verbatim}[fontsize=\footnotesize,frame=single,label={\small Python}]
import numpy as np
A = np.array([[10, 1, 1], [2, 10, 1], [2, 2, 10]],
             dtype=float)
b = np.array([12, 13, 14], dtype=float)
x = np.zeros(3)
for k in range(1, 4):
    for i in range(3):
        s = A[i] @ x - A[i, i] * x[i]
        x[i] = (b[i] - s) / A[i, i]
    err = np.abs(x - 1).max()
    print(k, x.round(5), f"{err:.1e}")
\end{Verbatim}
\begin{Verbatim}[fontsize=\footnotesize,frame=single,label={\small Output}]
1 [1.2   1.06  0.948] 2.0e-01
2 [0.9992  1.00536 0.99909] 5.4e-03
3 [0.99956 1.00018 1.00005] 4.4e-04
\end{Verbatim}
The same loop in Maple works with exact rational numbers, and
\texttt{evalf} converts the result to decimals at the end:
\begin{Verbatim}[fontsize=\footnotesize,frame=single,label={\small Maple}]
with(LinearAlgebra):
A := Matrix([[10, 1, 1], [2, 10, 1], [2, 2, 10]]):
b := Vector([12, 13, 14]):
x := Vector(3):
for k to 3 do
  for i to 3 do
    x[i] := (b[i] - add(A[i, j]*x[j], j = 1..3)
             + A[i, i]*x[i]) / A[i, i];
  end do;
end do:
evalf(x);        # 0.9995552, 1.000180160, 1.000052928
\end{Verbatim}
In Mathcad (with \texttt{ORIGIN} $=0$) the method is written as a small
program; the operator $\leftarrow$ is the local assignment of the
programming toolbar:
\begin{Verbatim}[fontsize=\footnotesize,frame=single,label={\small Mathcad}]
Seidel(A, b, x, m) :=
    | for k in 1..m
    |   for i in 0..rows(A)-1
    |     s <- sum(A_(i,j)*x_j, j) - A_(i,i)*x_i
    |     x_i <- (b_i - s)/A_(i,i)
    | return x
x0 := (0 0 0)^T
Seidel(A, b, x0, 3) = (0.99956 1.00018 1.00005)^T
\end{Verbatim}
Here \texttt{A} and \texttt{b} are entered as in the previous topic, and the
sum is entered with the summation operator over $j=0,1,2$.

\subsection{Exercises on the Seidel method}

\emph{1. Perform two iterations of the Seidel method for the system, starting from $x^{(0)}=(0;0;0)$.}

\bigskip
\begin{enumerate}[label=\arabic*., leftmargin=*]
\item $\begin{cases}7x_1+3x_3=27\\x_1-6x_2-2x_3=-13\\-2x_1-2x_2-10x_3=-30\end{cases}$
\item $\begin{cases}-10x_1-3x_2+x_3=20\\-2x_1+9x_2+2x_3=-23\\3x_1+8x_3=5\end{cases}$
\item $\begin{cases}-6x_1-x_2+2x_3=9\\-3x_1+8x_2-x_3=9\\-x_1+2x_2-8x_3=-13\end{cases}$
\item $\begin{cases}-11x_1-2x_2-3x_3=-11\\x_1+8x_2+3x_3=19\\-2x_1+2x_2-9x_3=22\end{cases}$
\item $\begin{cases}11x_1+3x_2-2x_3=-9\\-3x_1+10x_2-2x_3=-16\\-x_1-x_2-5x_3=-14\end{cases}$
\item $\begin{cases}9x_1-x_2+3x_3=9\\-2x_1+6x_2+x_3=3\\3x_1-8x_3=-24\end{cases}$
\item $\begin{cases}-12x_1-3x_2-3x_3=21\\-3x_1-9x_2+2x_3=8\\3x_2-8x_3=13\end{cases}$
\item $\begin{cases}11x_1+2x_2-3x_3=-6\\3x_1-11x_2+2x_3=4\\3x_1-3x_2-11x_3=-22\end{cases}$
\item $\begin{cases}-5x_1-x_2-x_3=4\\3x_1+10x_2-3x_3=-14\\-3x_2+7x_3=-8\end{cases}$
\item $\begin{cases}-9x_1-3x_2-3x_3=-12\\3x_1+11x_2-3x_3=-4\\-2x_1-3x_2-8x_3=24\end{cases}$
\item $\begin{cases}-9x_1+3x_2+x_3=-7\\3x_1-10x_2+x_3=-22\\3x_1+x_2-10x_3=-11\end{cases}$
\item $\begin{cases}-8x_1-2x_2-x_3=24\\-3x_1-8x_2+x_3=19\\-2x_1-2x_2+8x_3=24\end{cases}$
\item $\begin{cases}-8x_1+x_2+x_3=17\\-2x_1-8x_2=20\\-3x_1-2x_2+8x_3=34\end{cases}$
\item $\begin{cases}10x_1+2x_2+3x_3=-1\\-2x_1-8x_2-2x_3=12\\-x_1-x_2-8x_3=24\end{cases}$
\item $\begin{cases}9x_1-x_2+2x_3=-3\\x_1-5x_2-x_3=-15\\-2x_1-x_2+8x_3=-3\end{cases}$
\item $\begin{cases}5x_1=-5\\x_1+8x_2-2x_3=-11\\-3x_1-3x_2+10x_3=-21\end{cases}$
\item $\begin{cases}-8x_1+3x_2+x_3=18\\-2x_1-9x_2-x_3=-26\\x_1+2x_2-7x_3=-2\end{cases}$
\item $\begin{cases}7x_1+x_3=-19\\-3x_1-10x_2-3x_3=-7\\-2x_1+7x_3=20\end{cases}$
\item $\begin{cases}8x_1-3x_2-x_3=27\\2x_1-8x_2-2x_3=14\\2x_1+x_2-7x_3=5\end{cases}$
\item $\begin{cases}8x_1+x_2-3x_3=6\\6x_2-x_3=5\\-x_1-2x_2+6x_3=3\end{cases}$
\item $\begin{cases}7x_1+x_2-2x_3=-8\\2x_1+7x_2=10\\3x_1+2x_2-8x_3=14\end{cases}$
\item $\begin{cases}-7x_1+2x_2-x_3=-5\\-6x_2+2x_3=14\\-3x_1-2x_2+10x_3=14\end{cases}$
\item $\begin{cases}-8x_1+2x_2+3x_3=9\\x_1+7x_2+2x_3=-1\\3x_1+x_2-10x_3=25\end{cases}$
\item $\begin{cases}-7x_1+2x_2-2x_3=18\\4x_2-x_3=5\\2x_1-2x_2-9x_3=3\end{cases}$
\item $\begin{cases}-9x_1-3x_2+3x_3=-3\\-2x_1+7x_2-2x_3=7\\3x_1-2x_2-11x_3=-2\end{cases}$
\item $\begin{cases}6x_1-x_2-x_3=12\\2x_1-6x_2+x_3=-9\\-3x_3=-9\end{cases}$
\item $\begin{cases}9x_1-3x_2+x_3=-18\\3x_1+9x_2+2x_3=-21\\3x_1-9x_3=21\end{cases}$
\item $\begin{cases}-9x_1+3x_2+x_3=25\\-x_1-8x_2+2x_3=-26\\-x_1+2x_2+7x_3=-6\end{cases}$
\item $\begin{cases}-8x_1-2x_2=6\\8x_2+3x_3=-27\\x_1-3x_2-8x_3=17\end{cases}$
\item $\begin{cases}6x_1-2x_2=-20\\3x_1+10x_2+2x_3=3\\-x_1-2x_2+7x_3=8\end{cases}$
\end{enumerate}

\emph{2. Perform three iterations of the Seidel method for the $2\times2$ system, starting from $x^{(0)}=(0;0)$.}
\bigskip
\begin{multicols}{2}\raggedright
\begin{enumerate}[label=\arabic*., leftmargin=*]
\item $\begin{cases}8x_1+2x_2=-6\\-2x_1+5x_2=7\end{cases}$
\item $\begin{cases}8x_1=24\\-2x_1+8x_2=-14\end{cases}$
\item $\begin{cases}7x_1+2x_2=1\\x_1+6x_2=23\end{cases}$
\item $\begin{cases}7x_1-3x_2=13\\2x_1+4x_2=-6\end{cases}$
\item $\begin{cases}5x_1+3x_2=-13\\5x_2=-5\end{cases}$
\item $\begin{cases}6x_1-x_2=18\\-x_1+4x_2=-3\end{cases}$
\item $\begin{cases}6x_1+3x_2=-6\\-2x_1+4x_2=-18\end{cases}$
\item $\begin{cases}7x_1+x_2=22\\7x_2=7\end{cases}$
\item $\begin{cases}5x_1-2x_2=8\\3x_1+8x_2=-32\end{cases}$
\item $\begin{cases}7x_1=14\\3x_1+4x_2=14\end{cases}$
\item $\begin{cases}4x_1-2x_2=-16\\x_1+5x_2=-4\end{cases}$
\item $\begin{cases}6x_1-3x_2=-3\\2x_1+5x_2=-19\end{cases}$
\item $\begin{cases}5x_1-3x_2=4\\3x_1+6x_2=-21\end{cases}$
\item $\begin{cases}7x_1-3x_2=-13\\2x_1+8x_2=14\end{cases}$
\item $\begin{cases}5x_1-x_2=-1\\9x_2=9\end{cases}$
\item $\begin{cases}9x_1-3x_2=-21\\-3x_1+4x_2=10\end{cases}$
\item $\begin{cases}6x_1=-12\\5x_2=-15\end{cases}$
\item $\begin{cases}6x_1-3x_2=-12\\-x_1+6x_2=2\end{cases}$
\item $\begin{cases}6x_1-2x_2=14\\x_1+7x_2=17\end{cases}$
\item $\begin{cases}4x_1+3x_2=-16\\-3x_1+6x_2=12\end{cases}$
\item $\begin{cases}4x_1=-8\\2x_1+9x_2=-13\end{cases}$
\item $\begin{cases}9x_1-3x_2=-6\\-2x_1+9x_2=-32\end{cases}$
\item $\begin{cases}5x_1+3x_2=6\\-3x_1+8x_2=-33\end{cases}$
\item $\begin{cases}8x_1+x_2=4\\x_1+9x_2=-35\end{cases}$
\item $\begin{cases}5x_1-3x_2=-3\\-3x_1+5x_2=-11\end{cases}$
\item $\begin{cases}4x_1-x_2=7\\3x_1+7x_2=-18\end{cases}$
\item $\begin{cases}6x_1+3x_2=-33\\5x_2=-15\end{cases}$
\item $\begin{cases}9x_1=-36\\-3x_1+5x_2=12\end{cases}$
\item $\begin{cases}6x_1+3x_2=18\\3x_1+8x_2=35\end{cases}$
\item $\begin{cases}5x_1-2x_2=-3\\4x_2=-4\end{cases}$
\end{enumerate}
\end{multicols}

\emph{3. Applied problems (10 problems).}
\bigskip
\begin{enumerate}[label=\arabic*., leftmargin=*]
\item \textbf{(Comparing the methods.)} Perform three iterations of the Jacobi and of the Seidel method for the system $\begin{cases}7x_1-x_2=2\\3x_1-11x_2-3x_3=25\\3x_1-9x_3=9\end{cases}$, starting from $x^{(0)}=(0;0;0)$, and compare the errors $\|x^{(3)}-x^*\|_\infty$.
\item \textbf{(Comparing the methods.)} Perform three iterations of the Jacobi and of the Seidel method for the system $\begin{cases}-6x_1+3x_3=0\\2x_1-9x_2-2x_3=-16\\-x_1-2x_2+9x_3=-21\end{cases}$, starting from $x^{(0)}=(0;0;0)$, and compare the errors $\|x^{(3)}-x^*\|_\infty$.
\item \textbf{(Comparing the methods.)} Perform three iterations of the Jacobi and of the Seidel method for the system $\begin{cases}11x_1-3x_2+2x_3=-35\\2x_1-8x_2=-6\\-3x_1+2x_2+10x_3=-1\end{cases}$, starting from $x^{(0)}=(0;0;0)$, and compare the errors $\|x^{(3)}-x^*\|_\infty$.
\item \textbf{(Comparing the methods.)} Perform three iterations of the Jacobi and of the Seidel method for the system $\begin{cases}5x_1-x_2-x_3=-8\\-3x_1-11x_2+2x_3=15\\x_1+3x_2+7x_3=-12\end{cases}$, starting from $x^{(0)}=(0;0;0)$, and compare the errors $\|x^{(3)}-x^*\|_\infty$.
\item \textbf{(Comparing the methods.)} Perform three iterations of the Jacobi and of the Seidel method for the system $\begin{cases}7x_1+x_2+x_3=9\\2x_1-10x_2-2x_3=6\\-3x_1+3x_2+9x_3=21\end{cases}$, starting from $x^{(0)}=(0;0;0)$, and compare the errors $\|x^{(3)}-x^*\|_\infty$.
\item \textbf{(Comparing the methods.)} Perform three iterations of the Jacobi and of the Seidel method for the system $\begin{cases}10x_1+2x_2-2x_3=30\\x_1-7x_2+x_3=15\\2x_1-x_2+9x_3=-10\end{cases}$, starting from $x^{(0)}=(0;0;0)$, and compare the errors $\|x^{(3)}-x^*\|_\infty$.
\item \textbf{(Comparing the methods.)} Perform three iterations of the Jacobi and of the Seidel method for the system $\begin{cases}-6x_1-x_2+x_3=21\\x_1-8x_2-x_3=12\\x_1+3x_2+7x_3=-2\end{cases}$, starting from $x^{(0)}=(0;0;0)$, and compare the errors $\|x^{(3)}-x^*\|_\infty$.
\item \textbf{(Comparing the methods.)} Perform three iterations of the Jacobi and of the Seidel method for the system $\begin{cases}6x_1+x_3=1\\-x_1+6x_2-2x_3=-8\\-2x_1+3x_2+11x_3=8\end{cases}$, starting from $x^{(0)}=(0;0;0)$, and compare the errors $\|x^{(3)}-x^*\|_\infty$.
\item \textbf{(Comparing the methods.)} Perform three iterations of the Jacobi and of the Seidel method for the system $\begin{cases}-7x_1+2x_2=-20\\-3x_1+8x_2=-30\\2x_1+3x_2+11x_3=-16\end{cases}$, starting from $x^{(0)}=(0;0;0)$, and compare the errors $\|x^{(3)}-x^*\|_\infty$.
\item \textbf{(Comparing the methods.)} Perform three iterations of the Jacobi and of the Seidel method for the system $\begin{cases}-6x_1-x_2=13\\2x_1+8x_2+x_3=-14\\x_1+2x_2-6x_3=8\end{cases}$, starting from $x^{(0)}=(0;0;0)$, and compare the errors $\|x^{(3)}-x^*\|_\infty$.
\end{enumerate}

